\documentclass[11pt,a4paper]{article}
\usepackage[utf8]{inputenc}
\usepackage[T1]{fontenc}
\usepackage{lmodern}
\usepackage[margin=2.5cm]{geometry}
\usepackage{amsmath,amssymb}
\usepackage{booktabs}
\usepackage{longtable}
\usepackage{enumitem}
\usepackage{hyperref}
\usepackage{xcolor}
\usepackage{fancyhdr}
\usepackage{titlesec}
\usepackage{graphicx}
\usepackage{float}

\hypersetup{
    colorlinks=true,
    linkcolor=blue!60!black,
    citecolor=blue!60!black,
    urlcolor=blue!60!black
}

% Epistemic labels
\newcommand{\structural}{\textsc{[structural]}}
\newcommand{\functional}{\textsc{[functional]}}
\newcommand{\convergent}{\textsc{[convergent]}}
\newcommand{\exploratory}{\textsc{[exploratory]}}
\newcommand{\falsmark}{\textit{Falsified if}}

\title{\textbf{The Pharmacist's Cipher --- Companion Report}\\[0.5em]
\large A Morphological Framework, Visual-Textual Convergences,\\and Testable Predictions for MS\,408}

\author{Alessandro Placa\thanks{Corresponding author: \texttt{alessandro.placa@uroboro.io}}\\
\textit{Uroboro.io}}

\date{March 25, 2026\\[0.3em]
\small DOI: [to be assigned]\\[0.5em]
\small\textbf{Companion to:} Placa (2026), ``The Pharmacist's Cipher: Six Statistical Tests\\Supporting a Pharmaceutical Reading of the Voynich Manuscript (MS\,408)''\\
DOI: 10.5281/zenodo.19197846}

\begin{document}
\maketitle

\begin{center}
\small
\textbf{AI Disclosure:} Computational analysis and manuscript preparation were performed\\
with the assistance of Claude (Anthropic), an AI language model.\\
All hypotheses, interpretive decisions, and epistemic judgments are the author's.
\end{center}

\begin{abstract}
This companion report extends the statistical foundation established in Placa (2026)---hereafter ``Paper~I''---which presented six reproducible tests consistent with a pharmaceutical reading of the Voynich Manuscript (MS~408). Where Paper~I deliberately restricted itself to transcription-only analyses, this report presents the full working model that those statistics support: a proposed agglutinative morphology organized around solvent-based prefixes and therapeutic class markers, together with visual-textual convergence analyses using Yale IIIF high-resolution images.

The report serves four purposes: (1)~to describe the morphological framework in sufficient detail for independent evaluation, clearly separating structural facts from provisional glosses; (2)~to document convergences between textual profiles and manuscript illustrations, including 17~working botanical identifications produced through a ``sealed-envelope'' protocol; (3)~to present a record of 26+~falsified hypotheses, demonstrating methodological discipline; and (4)~to establish intellectual priority over 39~testable predictions at three confidence tiers, each with explicit falsification criteria.

This document is not a decipherment claim. It is a \textit{cartography of emerging implications}---a systematic inventory of the research directions that the framework opens, each recorded with its epistemic status, its anchorage to Paper~I, and its falsification path. Its purpose is to fix priority, demonstrate fertility, invite falsification, and prevent unattributed appropriation of research directions. For each implication, the report publishes the \textit{what} (the hypothesis) and the \textit{where it leads} (its consequences and testable predictions), while reserving for future work the full exposition of the specific convergence paths that generated each interpretive insight. Every section carries an explicit epistemic label (\structural{}, \functional{}, \convergent{}, or \exploratory{}) to ensure that no claim is presented at a higher confidence level than its evidence warrants.
\end{abstract}

\textbf{Keywords:} Voynich Manuscript, MS 408, agglutinative morphology, pharmaceutical codex, morphological framework, botanical identification, medieval pharmacopoeia, testable predictions

\newpage
\tableofcontents
\newpage

%---------------------------------------------------------------
\section*{How to Read This Report}
\addcontentsline{toc}{section}{How to Read This Report}

This document is not a second proof paper. It is an atlas of disciplined hypotheses generated by a framework that Paper~I has partially validated. Four principles govern its structure:

\begin{enumerate}[leftmargin=2em]
\item \textbf{This is not a decipherment claim.} The report proposes a working model, not a finished solution. Every section is explicitly marked for epistemic level.

\item \textbf{Claims are tiered by epistemic strength.} A five-level taxonomy is used throughout:
\begin{itemize}[leftmargin=1.5em]
\item \structural{} --- measurable distributional fact, reproducible from the corpus
\item \functional{} --- proposed role within the system, supported by distributional evidence and declared tests
\item \convergent{} --- interpretation supported by independent channels (text + image, or text + historical parallel)
\item \exploratory{} --- hypothesis with partial evidence and high potential yield
\item \textsc{[falsification criterion]} --- what would refute each claim
\end{itemize}

\item \textbf{Structural claims are sharply separated from glosses.} A segmentation claim (e.g., ``\textit{dy} is decomposable as \textit{d+y}'') is never presented at the same level as a semantic gloss (e.g., ``\textit{d} = medicine''). The former is distributional; the latter is interpretive.

\item \textbf{Many sections exist to establish priority and enable falsification,} not to assert conclusions. The report's value lies in making specific, testable, dated predictions---not in claiming they are already proven.

\item \textbf{This report is the second layer of a three-layer publication architecture.} Paper~I establishes the engine---the statistical evidence that a coherent system exists. This companion fixes the \textit{coordinates} of the implications that the engine generates: it publishes the \textit{what} (each hypothesis) and the \textit{where it leads} (its consequences and falsification path), while reserving for future dedicated work the full exposition of the specific convergence paths---visual, structural, codicological---that generated each interpretive insight. In other words: premises and landing points are public; the bridge between them is documented but not yet fully disclosed. This is not evasion: it is staged disclosure designed to fix intellectual priority while preserving the conditions for a structured program of verification and presentation.
\end{enumerate}

Every section and subsection carries one of these status labels. They are not decorative.

\subsection*{Relationship to Paper~I}

Paper~I (Placa, 2026) demonstrated six reproducible properties of the Voynich text consistent with a pharmaceutical register: cross-transcriber stability of morphological prefixes, a binary paradigm gap between \textit{s-} and \textit{sh-}, asymmetric lexical reuse between the Pharmaceutical and Stars sections, positional clustering of completion markers, section-specific morphological profiles, and a morpheme frequency hierarchy replicating pharmacological volume expectations. All six tests were computed from the transcription alone, with no dependence on illustrations or proposed glosses.

This companion report presents the morphological framework proposed to account for those patterns, together with visual-textual convergence analyses and a structured inventory of testable predictions. Nothing in this report duplicates Paper~I. Where a Paper~I result is referenced, it is cited by test number (e.g., ``Paper~I, Result~2'') and not re-derived.

%-------------------------------------------------------------[TRUNCATED FOR SUBMISSION — SEE FULL FILE IN GITHUB REPO]